\chapter{Variables and Singular Data Types}

A Python program does not begin with objects floating in isolation. It begins with names---entries in namespace dictionaries---that allow the runtime to reach heap-allocated objects, inspect them, combine them, and replace them during execution. For a C programmer, this is one of the first places where Python must be read differently: a Python variable is not a declared memory box with a fixed storage type. It is a name in a namespace, and that name can be bound to an integer object, a floating-point object, a boolean object, \texttt{None}, or another runtime object. The object carries the type through its \texttt{ob\_type} pointer; the name provides access to it.

This chapter studies how those bindings are established, how they interact with CPython's namespace architecture, and how the singular built-in value domains (\texttt{int}, \texttt{float}, \texttt{bool}, \texttt{None}) behave once bindings exist. The goal is not introductory syntax but runtime mechanics: keyword versus built-in masking, heap costs of reassignment on immutable types, arbitrary-precision integer representation, IEEE~754 floating-point boundaries, boolean inheritance from integer, and the algebraic precedence traps that arise when operators dispatch to object protocols rather than raw register arithmetic.

The chapter ends by connecting dynamic assignment with strong runtime operation checking. Python allows a name to move freely from one object to another, but it does not allow arbitrary operations between incompatible objects:

\begin{quote}
Names are flexible bindings. Objects keep their runtime type. Operations are checked against the objects involved.
\end{quote}
